\subsection{Abelian extensions and the abelian closure}

DEFINITION 1. --- An extension E of K is said to be abelian if it is Galois and its Galois group is commutative.

Since every subgroup of a commutative group is normal, Cor. 4 of V, p.~68 shows that every subextension of an abelian extension is abelian.

PROPOSITION 1. --- Let E be a Galois extension of K and Γ its Galois group. Let Δ be the derived group of Γ (I, p.~10, Def. 4) and F the field of invariants of Δ. For a subextension L to E to be abelian it is necessary and sufficient that it should be contained in F.

Firstly we note that F is also the field of invariants of the closure \(\bar{\Delta}\) of \(\Delta\) in \(\Gamma\) , and that \(\bar{\Delta}\) is a closed normal subgroup of \(\Gamma\) . By V, p.~68, Cor. 4, F is thus a Galois extension of K. Further, the Galois group of F over K is isomorphic to \(\Gamma/\bar{\Delta}\) and hence is commutative. Hence every subextension of F is abelian. Conversely let L be an abelian extension of K contained in E, and let \(\Pi\) be the Galois group of E over L. Since L is Galois, \(\Pi\) is a normal subgroup of \(\Gamma\) and the Galois group of L over K is isomorphic to \(\Gamma/\Pi\) (V, p.~68, Cor. 4). Therefore \(\Gamma/\Pi\) is commutative and \(\Pi\) contains \(\Delta\) , whence \(L \subset F\) .

COROLLARY. --- Let E be an extension of K and let \((\mathrm{E}_{i})_{i\in I}\) be a family of subextensions of E such that \(\mathrm{E}=\mathrm{K}\left(\bigcup_{i\in I}\mathrm{E}_{i}\right)\) . Suppose that each extension

\(E_{i}\) is abelian, then the same is true of E.

In the first place E is a Galois extension of K (V, p.~57, Prop. 1). If the field F is defined as in Prop. 1, then we have \(E_{i} \subset F\) for all \(i \in I\) , whence E = F.

An extension E of K is said to be an abelian closure of K if it is an abelian extension of K and every abelian extension of K is isomorphic to a subextension of E. Prop. 1 implies the existence of an abelian closure of K: for let \(K_{s}\) be a separable closure of K, with Galois group \(\Gamma\) and let \((\overline{\Gamma},\overline{\Gamma})\) be the closure of the derived group of \(\Gamma\) ; denote by \(K_{ab}\) the field of invariants of \((\overline{\Gamma},\overline{\Gamma})\) ; since every separable algebraic extension of K is isomorphic to a subextension of \(K_{s}\) (V, p.~45, Cor.). Prop. 1 shows that \(K_{ab}\) is an abelian closure of K. The Galois group of \(K_{ab}\) over K is canonically isomorphic to \(\Gamma/(\overline{\Gamma},\overline{\Gamma})\) . Let us now show the uniqueness of abelian closures: let E, \(E'\) be two abelian closures of K; by definition there exist K-homomorphisms \(u:E\to E'\) and \(v:E'\to E\) and Prop. 1 (V, p.~52) implies that \(v(u(E)) = E\) and \(u(v(E')) = E'\) , hence \(u\) is a K-isomorphism of \(E\) onto \(E'\) . Any other K-isomorphism of \(E\) onto \(E'\) is of the form \(u_1 = \sigma_0 \circ u\) with \(\sigma_0 \in \operatorname{Gal}(E'/K)\) ; since \(\operatorname{Gal}(E'/K)\) is commutative, the isomorphism \(\sigma \mapsto u \circ \sigma \circ u^{-1}\) of \(\operatorname{Gal}(E/K)\) onto \(\operatorname{Gal}(E'/K)\) is independent of \(u\) ; it is called the canonical isomorphism of \(\operatorname{Gal}(E/K)\) onto \(\operatorname{Gal}(E'/K)\) .

